\documentclass[11pt]{article}

\usepackage[a4paper,margin=28mm]{geometry}
\usepackage{amsmath,amssymb,mathtools}
\usepackage{booktabs}
\usepackage[hidelinks]{hyperref}
\usepackage[T1]{fontenc}
\usepackage{lmodern}

\newcommand{\KL}{D_{\mathrm{KL}}}
\newcommand{\Renyi}[1]{D_{#1}}
\newcommand{\Pnull}{P_0}
\newcommand{\Palt}{P_1}
\newcommand{\E}{\mathbb{E}}

\title{Technical Note: p-values, likelihood ratios, and information-theoretic divergence in FluxEMU}
\author{FluxEMU technical notes}
\date{6 September 2026}

\begin{document}
\maketitle

\section{Purpose}

Information-theoretic divergences are natural quantities for comparing isotope
observation laws, but they are not quantities that most experimentalists use in
routine interpretation. p-values, likelihood ratios, Type-I error, Type-II
error, and power are much more familiar. The aim of this note is to make the
relationship precise without treating a divergence as though it were itself a
p-value.

FluxEMU now distinguishes four related but different objects:
\begin{enumerate}
    \item a \emph{divergence} between complete observation laws;
    \item the \emph{likelihood-ratio evidence} in one realised data set;
    \item the \emph{p-value} obtained from the null tail of that evidence; and
    \item a finite-sample \emph{Type-II lower bound} at a declared Type-I budget.
\end{enumerate}

These quantities answer different scientific questions and should be reported
as such.

\section{Multinomial data and the KL deviance}

Suppose an experiment produces genuine isotopologue counts
\begin{equation}
K=(K_1,\ldots,K_m)\sim \operatorname{Multinomial}(n,p_0),
\end{equation}
where $p_0=(p_{01},\ldots,p_{0m})$ is a fully specified null MID. Let the
realised count vector be $k$, with empirical proportions
\begin{equation}
\widehat p_i=\frac{k_i}{n}.
\end{equation}

The multinomial log likelihood under a probability vector $p$ is, up to the
multinomial coefficient,
\begin{equation}
\ell(p;k)=\sum_{i=1}^m k_i\log p_i + C(k).
\end{equation}
The saturated multinomial model maximises this likelihood at
$p=\widehat p$. Therefore the likelihood-ratio deviance comparing $p_0$ with
the saturated model is
\begin{align}
G^2
&=2\{\ell(\widehat p;k)-\ell(p_0;k)\}\\
&=2\sum_{i=1}^m k_i\log\frac{\widehat p_i}{p_{0i}}\\
&=2n\KL(\widehat p\|p_0).
\end{align}
This identity is exact, with the usual convention $0\log 0=0$.

Thus KL divergence is already present inside the classical multinomial
likelihood-ratio statistic:
\begin{equation}
\boxed{G^2=2n\KL(\widehat p\|p_0).}
\end{equation}
Under the regular large-sample conditions of Wilks' theorem, $G^2$ converges to
a chi-square law. For a fully specified interior null on $m$ positive cells,
the usual saturated-model comparison has $m-1$ degrees of freedom. This yields
the familiar asymptotic p-value
\begin{equation}
p_{\mathrm{Wilks}}
\approx
\Pr\!\left\{\chi^2_{m-1}\ge 2n\KL(\widehat p\|p_0)\right\}.
\end{equation}
This is one direct bridge from KL divergence to a conventional p-value.
FluxEMU's exact simple-vs-simple p-value, developed below, does not require this
asymptotic approximation.

\section{Two fixed flux hypotheses}

For two fixed feasible flux states,
\begin{equation}
H_0:v=v_0,
\qquad
H_1:v=v_1,
\end{equation}
FluxEMU maps each flux state through stationary EMU prediction and the declared
observation model to two complete laws,
\begin{equation}
\Pnull=P(Y\mid v_0),
\qquad
\Palt=P(Y\mid v_1).
\end{equation}
The natural realised-data statistic is the log likelihood ratio
\begin{equation}
L(y)=\log\frac{\Palt(y)}{\Pnull(y)}.
\end{equation}
Large positive values favour the fixed alternative relative to the fixed null;
large negative values favour the null.

For a realised genuine-count observation $y_{\mathrm{obs}}$, define
\begin{equation}
\boxed{
p(y_{\mathrm{obs}})
=
\Pnull\!\left\{L(Y)\ge L(y_{\mathrm{obs}})\right\}.
}
\end{equation}
This is an exact one-sided likelihood-ratio p-value for the declared simple
null against the declared simple alternative. In words:

\begin{quote}
If $v_0$ were the true flux state, what is the probability of observing
likelihood-ratio evidence in favour of $v_1$ at least as strong as the evidence
actually observed?
\end{quote}

Because the statistic explicitly contains $P_1$, this p-value is
alternative-specific. Changing the alternative can change the ordering of
possible observations and therefore the p-value. It should not be presented as
a generic goodness-of-fit p-value for $v_0$ against all departures.

For a discrete count space, the tail probability defined with ``$\ge$'' is in
general super-uniform under the null. Consequently an exact discrete p-value
need not be continuously uniform under $H_0$ and may be conservative at a
nominal threshold.

\section{KL divergence as expected log evidence}

The same likelihood ratio gives KL divergence an operational interpretation.
If the alternative law is true,
\begin{align}
\E_{\Palt}[L(Y)]
&=\sum_y \Palt(y)\log\frac{\Palt(y)}{\Pnull(y)}\\
&=\KL(\Palt\|\Pnull).
\end{align}
If the null law is true,
\begin{align}
\E_{\Pnull}[L(Y)]
&=\sum_y \Pnull(y)\log\frac{\Palt(y)}{\Pnull(y)}\\
&=-\KL(\Pnull\|\Palt).
\end{align}
Therefore
\begin{equation}
\boxed{
\KL(\Palt\|\Pnull)
=\text{expected log evidence for }H_1\text{ when }H_1\text{ is true},
}
\end{equation}
while $\KL(\Pnull\|\Palt)$ is the magnitude of the expected log evidence in
the opposite direction under $H_0$.

This interpretation is often more useful experimentally than describing KL
divergence merely as a ``distance'' between distributions. KL divergence is
not symmetric and is not a metric; its direction has an operational meaning.

\section{Rényi divergence and likelihood-ratio tails}

For a finite order $\lambda>1$ and mutually absolutely continuous laws,
\begin{align}
\Renyi{\lambda}(\Palt\|\Pnull)
&=\frac{1}{\lambda-1}
\log\sum_y \Palt(y)^\lambda\Pnull(y)^{1-\lambda}\\
&=\frac{1}{\lambda-1}
\log\E_{\Pnull}\!\left[\exp\{\lambda L(Y)\}\right].
\end{align}
Hence Rényi divergence encodes an exponential moment of the same
likelihood-ratio evidence whose null tail defines the p-value. This provides a
direct conceptual route
\begin{equation}
\boxed{
\text{Rényi divergence}
\longrightarrow
\text{control of likelihood-ratio tails}
\longrightarrow
\text{finite-sample testing bounds}.
}
\end{equation}
The finite-sample Bruno--Vandenbroucque--Esposito result implemented in
FluxEMU uses both Rényi directions to lower-bound the optimal Type-II error
under a fixed Type-I constraint.

\section{Why the quantities should not be collapsed into one number}

Let
\begin{equation}
\alpha=P_0(\text{decide }H_1),
\qquad
\beta=P_1(\text{decide }H_0).
\end{equation}
The roles of the principal quantities are:

\begin{center}
\begin{tabular}{p{0.28\textwidth}p{0.64\textwidth}}
\toprule
Quantity & Scientific question \\
\midrule
$D_\lambda(P_1\|P_0)$ and $D_\lambda(P_0\|P_1)$
& How separated are the two complete observation laws? \\
$L(y)$
& Which fixed hypothesis does this realised data set favour, and by how much on the log-likelihood scale? \\
$p(y)=P_0\{L(Y)\ge L(y)\}$
& If the null were true, how often would evidence for this fixed alternative be at least this strong? \\
$\beta^*(\epsilon)$
& At Type-I budget $\alpha\le\epsilon$, how small can the best Type-II error be? \\
Finite-order Bruno lower bound
& What lower limit on $\beta^*(\epsilon)$ follows from the two Rényi divergences at the supplied order? \\
\bottomrule
\end{tabular}
\end{center}

A small p-value and a strong Type-II lower bound are not contradictory. The
first describes the realised observation under the null; the second describes
a limitation on the repeated-sampling discrimination problem at a fixed
Type-I budget. One observed data set can look unusual under $H_0$ even when the
experiment, considered before observing the data, has limited ability to
reliably separate $H_0$ from $H_1$.

\section{Data-versus-model and model-versus-model divergence}

FluxEMU uses divergences in two distinct settings that should not be conflated.

\subsection*{Data versus model}
For observed proportions $\widehat p$ and a predicted MID $p_v$,
\begin{equation}
\KL(\widehat p\|p_v)
\end{equation}
measures the discrepancy between the realised empirical distribution and a
model prediction. Under genuine multinomial counts, multiplying by $2n$ gives
the likelihood-ratio deviance against the saturated model.

\subsection*{Model versus model}
For two complete observation laws,
\begin{equation}
D_\lambda(P_1\|P_0),
\end{equation}
measures statistical separation between the two fixed flux hypotheses before a
particular count vector is observed. This is the quantity used by the
finite-sample Type-II lower bound.

\section{Large-deviation perspective}

In many independent-sample problems, probabilities of suitably defined rare
events have exponential form
\begin{equation}
P(\text{rare event})\asymp \exp\{-nI\},
\end{equation}
where $I$ is a large-deviation rate function. For empirical-distribution events
under an i.i.d. model, the relevant rate function is often a KL divergence.
This motivates the schematic relation
\begin{equation}
-\frac{1}{n}\log p \approx I,
\end{equation}
but this is not an exact general identity between a p-value and a divergence.
The event definition, test statistic, finite-sample discreteness, and competing
hypotheses all matter.

\section{Exact finite-sample implementation in FluxEMU}

For \texttt{likelihood\_ratio\_p\_value}, FluxEMU enumerates the complete
positive-probability null count space and sums
\begin{equation}
P_0(y)
\end{equation}
over all outcomes satisfying
\begin{equation}
L(y)\ge L(y_{\mathrm{obs}}).
\end{equation}
For explicitly independent observation blocks, the joint null space is the
Cartesian product of the blockwise count spaces and the joint log likelihood
ratio is additive.

The implementation preserves the following boundaries:
\begin{itemize}
    \item count totals must be genuine and explicitly declared;
    \item percentages, peak areas, normalised MIDs, and arbitrary intensities are not converted to pseudo-counts;
    \item exact structural zeros are retained;
    \item no pseudo-counts, smoothing, or support clipping are introduced;
    \item the exact enumeration has an explicit maximum-outcome limit;
    \item exceeding that limit raises an error rather than silently changing to a chi-square, Monte Carlo, or other approximation;
    \item the p-value is reported separately from the Rényi divergences and the Type-II lower bound.
\end{itemize}

\section{Worked FluxEMU mixer example}

Take
\begin{equation}
p_0=(0.25,0.75),
\qquad
p_1=(0.5,0.5),
\qquad
n=4,
\end{equation}
and suppose the realised genuine counts are
\begin{equation}
y_{\mathrm{obs}}=(3,1).
\end{equation}
The realised log likelihood ratio is
\begin{equation}
L(y_{\mathrm{obs}})
=\log\frac{P_1(3,1)}{P_0(3,1)}
\approx 1.6739764.
\end{equation}
Under $P_0$, evidence at least this strong occurs for counts $(3,1)$ and $(4,0)$,
so
\begin{align}
p(y_{\mathrm{obs}})
&=P_0(3,1)+P_0(4,0)\\
&=4(0.25)^3(0.75)+(0.25)^4\\
&=0.05078125.
\end{align}
At Rényi order $\lambda=2$, the complete $n=4$ multinomial laws satisfy
\begin{equation}
D_2(P_1\|P_0)\approx1.15072829,
\qquad
D_2(P_0\|P_1)\approx0.89257421.
\end{equation}
At Type-I budget $\epsilon=0.05$, the Bruno Type-II lower bound at that same
order is approximately
\begin{equation}
0.602476804.
\end{equation}
The p-value, the divergences, and the Type-II lower bound therefore describe
three different features of one simple testing problem; none should be used as
a synonym for another.

\section{Experimentalist-facing interpretation}

For a fixed comparison $v_0$ versus $v_1$, a useful report can present the
quantities together:

\begin{quote}
The observed isotope counts favour $v_1$ over $v_0$ with log likelihood ratio
$L(y)$. If $v_0$ were true, evidence at least this strong in favour of $v_1$
would occur with probability $p$. The complete observation laws have the
reported KL/Rényi divergences. At the chosen false-positive budget $\epsilon$,
the finite-sample Rényi result gives the reported lower bound on the optimal
false-negative probability.
\end{quote}

This keeps familiar inferential language while preserving the additional
information supplied by the information-theoretic formulation.

\begin{thebibliography}{9}

\bibitem{kl1951}
S. Kullback and R. A. Leibler,
``On Information and Sufficiency'',
\emph{Annals of Mathematical Statistics}, 22(1):79--86, 1951.

\bibitem{np1933}
J. Neyman and E. S. Pearson,
``On the Problem of the Most Efficient Tests of Statistical Hypotheses'',
\emph{Philosophical Transactions of the Royal Society A}, 231:289--337, 1933.

\bibitem{wilks1938}
S. S. Wilks,
``The Large-Sample Distribution of the Likelihood Ratio for Testing Composite Hypotheses'',
\emph{Annals of Mathematical Statistics}, 9(1):60--62, 1938.

\bibitem{bruno2026}
R. Bruno, A. Vandenbroucque, and A. R. Esposito,
``A Finite-Sample Strong Converse for Binary Hypothesis Testing via (Reverse) Rényi Divergence'',
arXiv:2601.09550v2, 2026.

\end{thebibliography}

\end{document}
